\section{A nonsplit extension of the 2D2 category}
\label{sec:nonsplit-2d2}

\begin{proposition}\label{prop:r37-model}
The ring $R8\_2\_37$ admits a unitary categorification.
\end{proposition}
\begin{proof}
We use the factor realization of a graded extension of the known
Haagerup--Izumi category for $C_4$ in
\cite[Section~3 and Theorem~4.1]{GISGraded}. It supplies automorphisms
$\alpha_j,\beta$ of a factor $M$, an endomorphism $\rho$, and unitaries
$u,v\in M$ such that
\[
\beta\alpha_j=\alpha_j\beta,\qquad
\beta\rho=\operatorname{Ad}(u)\alpha_1\rho\beta,\qquad
\beta^2=\operatorname{Ad}(v)\alpha_{1+z},\quad z\in\{0,2\}.
\]
The $\alpha_j$ form an honest $C_4$ action. The character identities
in \cite[Lemma~3.2, (3.10) and (3.18)]{GISGraded} give
$\alpha_2(u)=-u$ and $\alpha_2(v)=-v$.

Form $P=M\rtimes_{\alpha_2}C_2$, with implementing involution
$\lambda$. Grossman's construction
\cite[Section~4.1]{Grossman} extends the neutral generators by
\[
\widetilde\alpha_j(\lambda)=(-1)^j\lambda,
\qquad \widetilde\rho(\lambda)=\lambda.
\]
They generate the unitary $2D2$ category with simple sectors
$(1,g,r,gr)$, where $g=[\widetilde\alpha_1]$,
$r=[\widetilde\rho]$, and
\[
g^2=1,\qquad gr=rg,\qquad r^2=1+2r+2gr.
\]
This construction depends only on the algebra $1+\alpha_2$ in the
neutral category, so it applies to the neutral realization just chosen.

Extend $\beta$ to $b\in\operatorname{Aut}(P)$ by $b(\lambda)=\lambda$.
The factor identities extend to
\[
b\widetilde\rho
 =\operatorname{Ad}(u)\widetilde\alpha_1\widetilde\rho b,
\qquad
b^2=\operatorname{Ad}(v)\widetilde\alpha_{1+z}.
\]
Indeed, they hold on $M$ by construction. On $\lambda$ the two
right-hand sides are respectively $u(-\lambda)u^*=\lambda$ and
$v(-\lambda)v^*=\lambda$. The anticommutation of $u,v$ with
$\lambda$ is essential here.

Put $a=[b]$. These identities imply
\[
a^2=g,\qquad a^4=1,\qquad ara^{-1}=gr=a^2r,
\qquad ra=a^{-1}r.
\]
The full tensor category generated by $b$ and $\widetilde\rho$ has
exactly the eight simple sectors
\[
(1,a,a^2,a^3,r,ar,a^2r,a^3r).
\]
The four invertibles are distinct because $a^2=g\ne1$. If
$a^ir\cong a^jr$, Frobenius reciprocity puts $a^{j-i}$ in
$rr^*=1+2r+2gr$, so $i=j$. The noninvertibles have dimension
$2+\sqrt5$. The displayed relations show that these eight sectors
are closed under products and duals, and their associativity and
unitarity are inherited from $\operatorname{End}_0(P)$.

For $i,j\in\Z/4$, their multiplication is
\[
\begin{aligned}
a^ia^j&=a^{i+j},&
a^i(a^jr)&=a^{i+j}r,\\
(a^ir)a^j&=a^{i-j}r,&
(a^ir)(a^jr)&=a^{i-j}+2a^{i-j}r+2a^{i-j+2}r.
\end{aligned}
\]
The displayed basis maps to zero-based census labels
$(0,1,7,2,5,3,6,4)$. The exact verifier compares all $512$ structure
constants under this bijection. Thus it identifies an already
constructed category with $R8\_2\_37$; no coherent order-two
autoequivalence is assumed.
\end{proof}

\subsection{Odd isotypic multiplicities and involution algebras}

\begin{lemma}\label{lem:palindromic-fs}
Let $X$ be a self-dual simple in a pivotal fusion category over
$\C$. If a self-dual simple $S$ occurs with odd multiplicity in
$X^{\otimes n}$, then
\[
\nu_2(S)=\nu_2(X)^n.
\]
More generally, in a palindromic tensor word of self-dual simples,
every self-dual constituent of odd multiplicity has indicator equal
to the product of the indicators of the factors.
\end{lemma}
\begin{proof}
Use pivotal coherence to suppress associators and double-dual
identifications. Choose a self-duality $f:X\to X^*$ whose pivotal
transpose is $t(f)=\epsilon f$, where $\epsilon=\nu_2(X)\in\{1,-1\}$.
The reversed dual word has the same identical factors, so
\[
F=f^{\otimes n}:X^{\otimes n}\longrightarrow(X^{\otimes n})^*,
\qquad t(F)=\epsilon^nF.
\]
This construction needs no braiding: reversal changes neither the
object word nor the sequence of component morphisms. On the
$S$-isotypic block the nondegenerate form is $f_S\otimes B$, with
$t(f_S)=\nu_2(S)f_S$ and $B$ an invertible bilinear-form matrix on
the complex multiplicity space. Hence
\[
B^{\mathsf T}=\frac{\epsilon^n}{\nu_2(S)}B.
\]
An invertible alternating matrix has even size, which proves the
assertion for an odd multiplicity. Semisimplicity ensures that the
form on a self-dual isotypic block is nondegenerate. The same
calculation works for a palindrome, using the same chosen
self-duality at each pair of matching positions.
\end{proof}

\begin{proposition}\label{prop:r49-fs-morita}
The ring $R8\_1\_49$ has no categorification over $\C$.
\end{proposition}
\begin{proof}
Its invertibles are $G=\{0,3,4,5\}\cong C_2^2$, and its four
noninvertibles have dimension $d=2+\sqrt5$, with $d^2=1+4d$.
The published two-dimension theorem
\cite[Lemma~3.1]{SchNear} therefore replaces any hypothetical
categorification by a pseudo-unitary Galois conjugate with the
same fusion rules. Equip this conjugate with its canonical positive
spherical structure. This supplies the pivotal structure without
assuming one on the initial category.

Set $u=3$, $g=5$ and $\rho=6$. The source tensor gives
\[
\rho^*=\rho,\qquad
\rho^2=1+\sum_{h\in G}h\rho,\qquad
\rho^3=(1,4,4,1,1,1,5,4)
\]
in source-label order. The unit and every invertible occur once
in the cube. Lemma~\ref{lem:palindromic-fs} first gives
$\nu_2(\rho)=1$ and then $\nu_2(h)=1$ for every $h\in G$.
In the positive spherical structure of a pointed $C_2$ category,
the indicator of its nontrivial object is the sign of its
associator class. Thus the restriction to $H=\{1,u\}$ is
trivial, and its regular algebra $A=1\oplus u$ actually exists
and is separable. No central lift of $u$ is required.

Consider $\mathcal B={}_A\mathcal C_A$. The source relation
$u\rho u=g\rho$ gives two pointed $H$-by-$H$ orbits of size two,
$\{0,3\}$ and $\{4,5\}$, and the free noninvertible orbit
$\{1,2,6,7\}$. Each pointed orbit has stabilizer $C_2$ and
contributes two simple bimodules of dimension one, since every
complex projective representation of $C_2$ can be linearized.
The free orbit contributes the single simple $A\rho A$ of
bimodule dimension $4d/2=2d$. Thus $\mathcal B$ has four
invertibles, forming $G'\cong C_4$ or $C_2^2$, and one
noninvertible $Z$. Reciprocity and dimensions force
\[
Z^2=\sum_{h\in G'}h+8Z.
\]
For $G'=C_2^2$ this contradicts \cite[Theorem~1.1]{SchNear};
for $G'=C_4$ it contradicts the general cyclic-four criterion
\cite[Theorem~6.3]{R5} at $k=2$. The source-bound verifier
\texttt{verify\_r49\_fs\_morita.py} checks the cube, all double
orbits, and the same six-branch terminal certificate used below.
Hence the original ring cannot be categorified.
\end{proof}

\subsection{An inverse-Morita obstruction for the split involution}

\begin{proposition}\label{prop:r170-morita}
The ring $R8\_2\_170$ admits no fusion-category categorification
over $\C$.
\end{proposition}
\begin{proof}
Its source tensor has a faithful $C_2$ grading with neutral
labels $(0,5,6,7)$. Put $g=5$, $r=6$, and $u=3$. Its four
invertibles form $C_2^2$, with $u$ odd, and $uru=gr$.
The subgroup $H=\{1,u\}$ maps isomorphically to the grading
group. If a category with this tensor exists, multiply its
associator by the inflation of the inverse of its pointed
associator class on $H$. This is a normalized grading
three-cocycle, so it preserves the pentagon and every fusion
coefficient. After this twist $A=1\oplus u$ is a connected
separable algebra. This step does not assume that $A$ exists
before twisting.

Consider the fusion category $\mathcal B={}_A\mathcal C_A$.
For a free bimodule $F(Y)=A\otimes Y\otimes A$, adjunction gives
\[
\dim\operatorname{End}_{\mathcal B}(F(Y))
 =\dim\operatorname{Hom}_{\mathcal C}(Y,A\otimes Y\otimes A)
 =|\{(h,k)\in H^2:hYk\cong Y\}|.
\]
The two pointed double orbits are $\{0,3\}$ and $\{4,5\}$,
each of stabilizer size two. Each free bimodule therefore
splits into two distinct simple bimodules, once each: its
semisimple endomorphism algebra has dimension two. The
underlying support of either summand is stable under $H$
on both sides, so each contains both orbit objects once.
Dividing its underlying FP dimension by $\FPdim(A)=2$
shows that it is invertible. Thus these orbits give four
invertibles in $\mathcal B$.

The four noninvertibles form one double orbit
$\{1,2,6,7\}$ with trivial stabilizer. Hence $F(r)$ is a
single simple $Z$ of dimension $2(2+\sqrt5)=4+2\sqrt5$.
Every simple bimodule occurs in a free one, so this is the
complete list. The four invertibles form a group $G$ equal
to $C_4$ or $C_2^2$. The unique noninvertible is self-dual
and fixed by $G$; reciprocity and dimensions give
\[
Z^2=\sum_{h\in G}h+8Z.
\]
For $G=C_2^2$, this contradicts
\cite[Theorem~1.1]{SchNear}. For $G=C_4$, it contradicts
the general cyclic-four central-support criterion of
\cite[Theorem~6.3]{R5}, at parameter $k=2$ and moment
height six. The archived optimizer-free replay checks all
six branches: two force a fusion coefficient $1/2$, and
four have exact separating vectors with nonnegative
$A^{\mathsf T}y$ and respective values
$b^{\mathsf T}y=-800,-220,-320,-80$. This terminal
criterion applies to arbitrary complex categorifications
through its proved Galois pseudo-unitarity reduction.
Both possible groups are excluded. No classification of
the neutral $2D2$ category is used.
\end{proof}

\section{The cyclic order-six Haagerup--Izumi algebra dual}
\label{sec:hi6-dual}

\begin{proposition}\label{prop:r535-model}
The dual of the algebra $1+\rho$ in a known unitary Haagerup--Izumi
category for $C_6$ has Grothendieck ring $R8\_2\_535$.
\end{proposition}
\begin{proof}
Set $d=3+\sqrt{10}$ and $D=6(1+d^2)=120+36\sqrt{10}$.
The category and its algebra $1+\rho$ exist by
\cite[Example~4.5 and condition~(2.24)]{GICenters}. Denote the
unitary dual by $\mathcal E$. Izumi's description
\cite[Proposition~7.4 and Theorem~7.5]{Izumi3n} gives its simple
objects and dimensions as
\[
\begin{array}{c|rrrrrrrr}
&1&b&r&rb&e_1&e_2&s_1&s_2\\\hline
\FPdim&1&1&d&d&d+1&d+1&d-1&d-1.
\end{array}
\]
Here $b=\beta_3$, $e_j=\eta_j$, and $s_1,s_2$ are the two
$\sigma$ objects. Its canonical module has seven simples
$(i_0,\ldots,i_5,\tau)$ with relative dimensions
$(1,\ldots,1,d-1)$.

Let $P$ be the cyclic permutation matrix on the six $i$-labels.
The upper-left six-by-six blocks of the eight action matrices are
\[
I,\quad P^3,\quad I,\quad P^3,\quad
P+P^{-1},\quad P^2+P^{-2},\quad0,\quad0.
\]
The matrices of $1,b$ have zero off-diagonal $\tau$ entries and
$\tau$ diagonal $1$. Each of the other six has every off-diagonal
$\tau$ entry equal to $1$, with respective $\tau$ diagonals
$(5,5,6,6,4,4)$. These formulas follow from the decompositions in
\cite[Proposition~7.4]{Izumi3n}; the last entries also follow from
the relative FP dimension vector.

These matrices generate a commutative algebra of dimension five.
Their three independent dependencies are
\[
M_1-M_b-M_r+M_{rb}=0,\qquad
M_1-M_r+M_{s_1}=0,\qquad M_1-M_r+M_{s_2}=0.
\]
The complex fusion algebra is semisimple. A noncommutative simple
block would have dimension at least four and would have to lie in
the three-dimensional kernel of this representation. Hence the
entire fusion algebra of $\mathcal E$ is commutative.

Write $d'=3-\sqrt{10}$ and $D'=120-36\sqrt{10}$. Fourier
diagonalization of the module yields five characters:
\[
\begin{array}{c|rrrrrrrr|c}
&1&b&r&rb&e_1&e_2&s_1&s_2&f\\\hline
\chi_{\rm FP}&1&1&d&d&d+1&d+1&d-1&d-1&D\\
\chi'_{\rm FP}&1&1&d'&d'&d'+1&d'+1&d'-1&d'-1&D'\\
u&1&-1&1&-1&1&-1&0&0&6\\
v&1&1&1&1&-1&-1&0&0&6\\
w&1&-1&1&-1&-2&2&0&0&12.
\end{array}
\]
There are three further one-dimensional characters. Orthogonality
to the displayed rows forces each to have the form
\begin{equation}\label{eq:hi6-missing}
(1,\epsilon,-1,-\epsilon,0,0,s,t),\qquad
\epsilon\in\{1,-1\},\quad s+t=1+\epsilon.
\end{equation}

We determine their codegrees using the center of the known category,
which is Morita invariant. The complete published inventory
\cite[Lemma~3.5 and Theorem~3.9]{GICenters} is
\[
\begin{array}{c|rrrrrr}
\text{simple dimension}&1&d^2&1+3d&2+6d&3d&6d\\\hline
\text{multiplicity}&2&2&4&16&4&18.
\end{array}
\]
The dimensions squared sum to $D^2$. Formal-codegree induction
\cite{OstrikCode} associates to each character a simple constituent
of $I(1)$ of dimension $D/f$. Thus the possible codegrees are
\[
D,\quad D',\quad12,\quad6,\quad4\sqrt{10},\quad2\sqrt{10}.
\]
Formal codegrees are totally positive: in this commutative case each
is the squared norm of a character and every Galois conjugate is
another formal codegree. This eliminates the last two values.
For a missing character, the Perron bound gives
\[
f=4+|s|^2+|t|^2
 \leq4+2(d-1)^2=8+8d<D.
\]
So no missing codegree is $D$. The five known rows form a
Galois-stable set; a missing codegree $D'$ would have a missing
Galois conjugate of codegree $D$, also impossible. Finally
$\sum_\chi f_\chi^{-1}=1$ and $D^{-1}+(D')^{-1}=1/6$ force the
three missing codegrees to be $6,6,12$.

The object $b$ permutes $s_1,s_2$. If it fixed both, a missing
character with $\epsilon=-1$ would have $s=t=0$ and codegree $4$.
All three missing rows would then have $\epsilon=1$, but the vectors
in \eqref{eq:hi6-missing} with $\epsilon=1$ span a space of dimension
at most two. Therefore $b$ exchanges $s_1$ and $s_2$, and
$t=\epsilon s$. There is exactly one missing row with $\epsilon=1$,
namely $(1,1,-1,-1,0,0,1,1)$ of codegree $6$. The two remaining
rows have $\epsilon=-1$, one of each codegree $6,12$. Complex
conjugation preserves the missing set, $\epsilon$, and the
codegree, so each of these two rows is real. Their norms and
orthogonality determine, up to exchanging $s_1,s_2$, the full list
\[
\begin{array}{c|rrrrrrrr|c}
x&1&1&-1&-1&0&0&1&1&6\\
y&1&-1&-1&1&0&0&1&-1&6\\
z&1&-1&-1&1&0&0&-2&2&12.
\end{array}
\]
All characters are real, hence all simple objects are self-dual.
Character inversion now determines every structure constant:
\[
N_{ij}^{k}=\sum_\chi
 \frac{\chi(i)\chi(j)\chi(k)}{f_\chi}.
\]
The resulting tensor is identified with $R8\_2\_535$ by the
basis map $(0,1,3,6,4,5,2,7)$. The verifier reconstructs all $512$
coefficients exactly and checks the canonical module equations.
No multiplicity bound or completeness claim about the census is
used to determine the unknown dual ring.
\end{proof}

\section{Algebraic rigidity of the 2D2 module graph}
\label{sec:algebraic-2d2}

The next argument gives four general complex obstructions. It does not assert that every
categorification of the rank-four $2D2$ ring is unitary: the specific
module graph is part of the hypothesis.

\begin{lemma}\label{lem:algebraic-2d2}
Let $\mathcal D$ be a pseudo-unitary fusion category and
$\mathcal M$ an indecomposable module category with a simple object
$X$. Suppose its rooted principal graph is $2D2$, with modulus
$\delta=\sqrt{3+\sqrt5}$. Then the associated algebraic shaded
planar algebra is isomorphic to the known unitary $2D2$ planar
algebra. In particular its dual even category has rank six.
\end{lemma}
\begin{proof}
We first justify the algebraic planar-algebra construction.
Schaumann's module-trace theorem
\cite[Proposition~5.8]{SchaumannTraces} supplies positive FP trace
weights on $\mathcal M$. His Theorem~4.5 constructs natural
double-adjoint isomorphisms for module functors between any two
traced module categories, compatible with composition. Applied to
$\mathcal D$ and $\mathcal M$, it equips the full Morita linking
category with a pivotal structure. Corollary~4.6 and
Proposition~5.10(iv) give the compatible spherical dual. Normalize
the module trace so that the two loop values of $X$ are $\delta$.

For completeness, the graph embedding in this setting requires no
star operation. Act on a column of the linking category, and
decompose every object into simples. A morphism acts on the
resulting path spaces as a direct sum of matrices, faithfully.
Choose dual bases for the nondegenerate trace pairings. Composition
is matrix multiplication, while evaluation and coevaluation are
the weighted graph cups and caps. Thus all planar operations are
preserved. This is the algebraic pivotal construction of
\cite[Theorems~3.45 and~3.70]{EH}, with the cup and cap calculation
of its Theorem~3.75 performed with dual bases. Its dagger-specific
Theorem~3.77 is not needed. Denote the resulting planar algebra by
$\mathcal P$ and its embedding into the graph planar algebra by
$\mathcal P\subseteq\mathcal G$.

The graph is two-supertransitive. Its depth-three Jones--Wenzl
corner is $f_3\mathcal P_{3,+}f_3=\C p\oplus\C q$, where
$p,q$ are orthogonal primitive idempotents and $p+q=f_3$. The two
depth-three PF weights coincide, so $T=p-q$ has trace zero and
$T^2=f_3$. The bilinear trace pairing is nondegenerate, and
\[
\mathcal P_{3,+}=\mathrm{TL}_{3,+}\oplus\C T
\]
is a trace-orthogonal decomposition. Pairing a cap of $T$ with
every element of $\mathcal P_{2,+}$ shows that it vanishes, since
the corresponding cups lie in $\mathrm{TL}_{3,+}$. Rotation
preserves the one-dimensional orthogonal complement, so $T$ is
an uncappable rotational eigenvector, whose eigenvalue is a cube
root of unity. We verify the exclusion of the two primitive
roots directly below; this also independently replays
\cite[Lemma~3.2]{MorrisonPenneys2D2} without a self-adjointness
hypothesis.

We describe the exact complex calculation needed at this point.
Label the graph vertices $0,\ldots,7$, with edges
\[
01,\ 12,\ 23,\ 24,\ 35,\ 45,\ 56,\ 67,
\]
rooted at $0$ and with even vertices $0,2,5,7$. Put
$s=\sqrt5$, $\varphi=(1+s)/2$, and assign weights
\[
(w_0,\ldots,w_7)=(1,1,2+s,\varphi,\varphi,2+s,1,1).
\]
These are the PF weights with odd entries divided by $\delta$.
For a path $p=(p_0,p_1,p_2,p_3)$, set
$A_p=\sqrt{w_{p_1}w_{p_2}}$. The rectangular lopsided matrix is
$T_{\rm rect}=A T_{\rm sph}A^{-1}$. Thus the coefficient at a
loop $(v_0,\ldots,v_5)$ is multiplied by
$\sqrt{w_{v_1}w_{v_2}/(w_{v_5}w_{v_4})}$.
In these coordinates a Temperley--Lieb projection has entries
$w_{p_i}/\sum_{a\sim p_{i-1}}w_a$ on paths differing only
at the backtracking vertex $p_i$, and zero otherwise. Caps
therefore give unweighted sum equations. Two-click rotation
replaces $(v_0,\ldots,v_5)$ by $(v_4,v_5,v_0,v_1,v_2,v_3)$,
with coefficient $w_{v_0}w_{v_5}/(w_{v_2}w_{v_1})$.

For a primitive root $\omega$, work in
$\Q(s,\omega)$, $\omega^2+\omega+1=0$. The cap and
$\omega$-rotation equations on the $152$ loops have rank
$144$. After their elimination, let $E_j$ be the $90$
distinct normalized nonzero entries of $T^2-f_3$, in the
deterministic path-matrix order of the verifier. Direct
polynomial arithmetic gives
\[
\begin{split}
1={}&(-11-5s)E_{26}
 +\left(-25-11s+\frac{5+3s}{2}\omega\right)E_0\\
 &+\left(-12-6s+\frac{9+3s}{2}\omega\right)E_1
 +(44+20s)E_4\\
 &+\bigl(4+2s-(7+3s)\omega\bigr)E_{13}.
\end{split}
\]
The portable checker reconstructs the linear elimination and
every $E_j$ from the graph before checking this identity. It
excludes both primitive cubic roots simultaneously, without
complex conjugation equations. Hence the eigenvalue is one.

For eigenvalue one the cap and rotation system again has rank
$144$. Its eight free loop
coordinates $x_0,\ldots,x_7$ may be chosen, in order, at
\[
012321,\ 212323,\ 212354,\ 212453,\
232323,\ 232354,\ 232453,\ 235453.
\]
Put $k=(3-s)/2$ and $c=(5s-11)/2>0$.
On this eight-dimensional linear space, the ideal of all
path-matrix entries of $T^2-f_3$ is generated by
\begin{equation}\label{eq:2d2-complex-ideal}
\begin{gathered}
x_0+\frac{3+s}{2}x_7,\qquad
x_1+\frac{3-s}{2}x_7,\qquad
x_2+2x_5,\qquad x_3+2x_6,\\
x_4+\frac{s-1}{2}x_7,\qquad
x_5x_6-c,\qquad x_7^2-k^2.
\end{gathered}
\end{equation}
This is an equality of polynomial ideals over $\Q(\sqrt5)$,
not just equality of their real or self-adjoint solution sets.
The portable certificate generates all square equations from
the graph and verifies both inclusions: reduction proves that
\eqref{eq:2d2-complex-ideal} implies every square equation, and
explicit polynomial identities of degree at most three express
its seven generators in the original equations. The linear rank
and all coefficients are checked exactly.

Thus every complex solution is $\pm T(\lambda)$ for
$\lambda\in\C^\times$, with free coordinates
\[
\left(-1,-k^2,-2\sqrt c\lambda,-2\sqrt c\lambda^{-1},
 -\frac{s-1}{2}k,\sqrt c\lambda,
 \sqrt c\lambda^{-1},k\right).
\]
The base solution $T(1)$ is self-adjoint after the diagonal
change to spherical coordinates. Assign an oriented-edge
character $\lambda$ to the edge $3\to5$, its inverse to the
reverse edge, and $1$ to all others. Multiplying each loop
coefficient by its edge-character product is a graph planar
algebra automorphism: the factors on internal edges cancel in
every state sum. The eight coordinates have winding numbers
$(0,0,1,-1,0,1,-1,0)$, so this gauge carries $T(1)$ to
$T(\lambda)$. The two signs generate the same subalgebra.

After applying this gauge, $\mathcal P$ contains the planar
subalgebra generated by $T(1)$. By
\cite[Lemma~3.4 and Theorems~7.8--7.9]{MorrisonPenneys2D2},
this is the known unitary $2D2$ planar algebra. Its positive
box dimensions are the rooted-loop dimensions of $2D2$,
exactly those of $\mathcal P$. The containment is therefore
equality in every positive box space; one-click rotation gives
equality in the negative spaces of positive degree, and both
zero-box spaces are $\C$. This rules out additional
higher-depth data without using positivity on $\mathcal P$.
The known dual graph has level sizes $(1,1,1,2,4,1)$, hence
six even vertices.
\end{proof}

\begin{proposition}\label{prop:2d2-four-general}
The rings $R8\_2\_152$, $R8\_2\_153$, $R8\_2\_166$ and
$R8\_2\_495$ admit no fusion-category categorification over $\C$.
\end{proposition}
\begin{proof}
Each tensor has a faithful $C_2$ grading with rank-four neutral
ring $g^2=1$, $gr=rg$, $r^2=1+2r+2gr$. Its odd component
contains $X$ with $XX^*=X^*X=1+r$. The source witnesses are
\[
\begin{array}{c|c|c|c}
\text{ring}&(1,g,r,gr)&X&
\text{graph vertices }0,\ldots,7\\\hline
R8\_2\_152&(0,6,5,7)&1&(0,1,5,3,4,7,2,6)\\
R8\_2\_153&(0,6,5,7)&1&(0,1,5,3,4,7,2,6)\\
R8\_2\_166&(0,5,7,6)&3&(0,3,7,1,2,6,4,5)\\
R8\_2\_495&(0,1,6,3)&2&(0,2,6,4,5,3,7,1).
\end{array}
\]
Both left and right incidence matrices for tensoring with $X$
are the rooted graph of Lemma~\ref{lem:algebraic-2d2}. The object
$r$ generates the full neutral category, so both even parts
have rank four.

Suppose a complex categorification exists. Its neutral category
has exactly two FP dimensions. By \cite[Lemma~3.1]{SchNear},
it is Galois conjugate to a pseudo-unitary category. Apply the
same field automorphism to the entire category and its odd
module; their based graphs are unchanged. Lemma~\ref{lem:algebraic-2d2}
then forces the dual even part to have rank six, contradicting
the rank-four graph computed from the tensor. Equivalently,
the odd component is an invertible neutral bimodule and its
dual category is the rank-four neutral category. No pivotal
structure on the original rank-eight category is assumed.
\end{proof}

\section{Pointed support in a twisted half-braiding}
\label{sec:twisted-pointed-support}

\begin{proposition}\label{prop:twisted-two}
The rings $R8\_2\_314$ and $R8\_2\_432$ admit no
fusion-category categorification over $\C$.
\end{proposition}
\begin{proof}
Both tensors have a fusion subring $\mathcal D$ with basis
$(1,g,r,gr)$ and rules
\[
g^2=1,\qquad gr=rg,\qquad r^2=g^\epsilon+r+gr,
\]
where $\epsilon=1$ for $R8\_2\_314$ and $\epsilon=0$ for
$R8\_2\_432$. The complement $\mathcal M$ is stable under
both $\mathcal D$ actions. Its four simples are $aX$ for
$a\in\Irr\mathcal D$, with
\[
\begin{gathered}
Xa\cong F(a)X,\qquad X^2=g^\epsilon+X+gX,\\
F(1)=1,\quad F(g)=g,\quad F(r)=gr,\quad F(gr)=r.
\end{gathered}
\]
The complete source witnesses are
\[
\begin{array}{c|c|c|c}
\text{ring}&(1,g,r,gr)&X&(X,gX,rX,grX)\\\hline
R8\_2\_314&(0,7,1,2)&3&(3,4,6,5)\\
R8\_2\_432&(0,3,4,5)&6&(6,7,1,2).
\end{array}
\]
Here $F$ is initially only notation for the displayed label
permutation. We explain why a hypothetical categorification
forces the coherence needed below.

The functor $L:\mathcal D\to\mathcal M$, $a\mapsto aX$, is
an equivalence of left module categories: it maps simples
bijectively to simples and is fully faithful on their Hom
spaces. Transport the right $\mathcal D$ action through $L$.
Left module endofunctors of the regular category are right
tensor functors, so this action supplies a strong monoidal
functor $F:\mathcal D\to\mathcal D$ with coherent natural
isomorphisms $Xa\cong F(a)X$. Its displayed permutation of
simples makes it a tensor autoequivalence.

Applying this relation twice gives coherent isomorphisms
$X^2a\cong F^2(a)X^2$. The decomposition
$\mathcal C=\mathcal D\oplus\mathcal M$ is a decomposition
of $\mathcal D$-bimodule categories, so its projection to
$\mathcal M$ preserves these constraints. Write $z=1\oplus g$.
The $\mathcal M$ part of $X^2$ is $zX$. Using $Xa\cong F(a)X$
and cancelling the fully faithful functor $L$ gives
\[
\eta_a:zF(a)\longrightarrow F^2(a)z.
\]
If $J^F,J^{F^2}$ denote the tensorators, its compatibility is
\[
\eta_{ab}(\mathrm{id}_z\otimes J^F_{a,b})
 =(J^{F^2}_{a,b}\otimes\mathrm{id}_z)
  (\mathrm{id}_{F^2(a)}\otimes\eta_b)
  (\eta_a\otimes\mathrm{id}_{F(b)}),
\]
with associators understood. Equivalently this is the bimodule
constraint of the $\mathcal M$ projection of multiplication
$\mathcal M\boxtimes_{\mathcal D}\mathcal M\to\mathcal C$.
Applying a monoidal quasi-inverse of $F$, and using
$F^{-1}(z)\cong z$, gives an $F$-twisted half-braiding
\begin{equation}\label{eq:twisted-half-braiding}
\theta_a:za\longrightarrow F(a)z,\qquad
\theta_{ab}=(J^F_{a,b}\otimes\mathrm{id}_z)
 (\mathrm{id}_{F(a)}\otimes\theta_b)
 (\theta_a\otimes\mathrm{id}_b).
\end{equation}
No identification of $F^2$ with the identity is required.

Fix the two summands $z_0=1,z_1=g$. Since $r\ne gr$ and
$F(r)=gr$, the map $\theta_r$ must send $z_i r$ to
$F(r)z_{i+1}$, with indices modulo two. Thus it exchanges the
two support labels. Equation~\eqref{eq:twisted-half-braiding}
for $a=b=r$ applies two exchanges, and hence $\theta_{r\otimes r}$
preserves these labels. Its associators and tensorator preserve
the relevant direct-sum idempotents by naturality.

Choose a nonzero inclusion $v:r\to r\otimes r$, which exists
in both rings. Naturality requires
\[
\theta_{r\otimes r}(\mathrm{id}_z\otimes v)
 =(F(v)\otimes\mathrm{id}_z)\theta_r.
\]
On $z_i r$ the left side lands in $F(r\otimes r)z_i$ and
the right side in $F(r\otimes r)z_{i+1}$. Both maps are
nonzero: $v,F(v)$ are split monomorphisms, tensoring preserves
them, and the $\theta$ maps are isomorphisms. The two target
direct summands have zero intersection, a contradiction.
The argument uses only semisimplicity, rigidity and tensor
coherence, with no pivotality assumption.
\end{proof}

\section{A forced quantum-group Morita dual}
\label{sec:r342-morita}

\begin{proposition}\label{prop:r342-morita}
The ring $R8\_2\_342$ has no categorification over $\mathbb C$.
\end{proposition}
\begin{proof}
Suppose $\mathcal C$ categorifies this ring. The invertible
$g=5$ fixes simples $3,7$. Its pointed associator is therefore
trivializable, since either fixed simple gives a rank-one
module category over $\langle g\rangle$. Thus $A=1\oplus g$
is a connected separable algebra. Put $\mathcal M=\mathcal C_A$
and let $\mathcal E$ be its Morita dual.

The simple module labels are
\[
(0,0),(1,0),(3,+),(3,-),(4,0),(7,+),(7,-).
\]
The exact certificate \texttt{verify\_r342\_inverse\_morita.py}
exhausts the two-dimensional signed representations on the
two split labels. Certified annihilating polynomials and the
three source generators $1,4,6$ leave eight generator tuples
and exactly one full representation up to sign gauge.
Consequently all seven-by-seven left action matrices $L_i$
are forced. In particular $L_g=I_7$.

The distinguished induction Gram matrix is $2I_7$. Its
nonnegative integral columns must be exactly two copies of
each standard basis vector. Hence $\mathcal E$ has fourteen
simples, two above each module label. Their right action
matrices $R_y$ commute with every $L_i$ and satisfy
\[
Q_a:=\sum_i(L_i)_{aa}L_i
 =\sum_yR_y[:,a]R_y[:,a]^t,
\qquad
0\le(R_y)_{ba}\le\lfloor\sqrt{(Q_a)_{bb}}\rfloor.
\]
These are linking-category associativity identities.
The common commutant has dimension eleven; a prescribed
standard first column leaves four free matrix entries.
Exhausting the displayed bounds gives candidate counts
$(2,4,2,2,2,2,2)$. The first domain contains the identity
and a unique nontrivial involutive permutation $P$, representing
the second invertible $u$. Right multiplication by $P$ pairs
every domain freely. Of the resulting two joint candidate
tuples, exactly one satisfies all seven Gram identities.

The certificate reconstructs the full dual fusion tensor,
rather than stopping at these necessary identities. For each
product its coefficient vector $n$ obeys
\[
\sum_wn_wR_w=R_yR_z,\qquad
n_1=\delta_{z,y^*},\qquad n_u=\delta_{zu,y^*}.
\]
The first column bounds each coefficient by the corresponding
entry of $R_yR_z$. Complete enumeration forces the ten fusion
matrices with labels $0,1,4,5,6,7,10,11,12,13$.
The fourteen words
\[
1,u,4,5,6,7,10,11,12,13,44,46,4\,11,4\,13
\]
have independent class vectors. Their associativity therefore
forces all remaining matrices. The resulting tensor passes
all original product domains and all $14^4$ associativity
equations.

The labels $(0,2,5,6,9,10,13)$ form a fusion subcategory
$\mathcal B$, with complement $\mathcal Bu$ and faithful
$C_2$ grading. The restriction $\mathcal M|_{\mathcal B}$
is regular. Its based ring is $\mathrm{PSU}(3)_5$:
the generator matrix in \cite[Section~2.4]{BRWRealizing}
matches under the local map $(0,4,5,6,1,3,2)$, and its
seven powers on the unit span the whole ring.
Conjugation by $u$ induces precisely its duality permutation
$(0,1,3,2,4,6,5)$. The checker also verifies
$\mathcal B_{\mathrm{ad}}=\mathcal B$ and that its based
automorphism group is $C_2$.

The Kazhdan--Wenzl theorem, in the general monoidal form
\cite[Theorem~5.1]{FMRReconstruction}, applies to
$\mathcal B\boxtimes\operatorname{Vec}_{C_3}$, whose based
ring is $\mathrm{SU}(3)_5$. Restrict its resulting equivalence
to adjoint subcategories. The grading three-cocycle is trivial
there, so $\mathcal B$ is a Galois conjugate of the standard
quantum $\mathrm{PSU}(3)_5$ category. This step assumes
neither a braiding nor pivotality.

The gauge-autoequivalence kernel of $\mathrm{SU}(3)_5$ is
trivial by \cite[Theorem~1.1 and Section~3]{EdieSU3}.
Since the standard category factors as
$\mathrm{PSU}(3)_5\boxtimes\mathcal P$, with $\mathcal P$
pointed of rank three, any based-trivial tensor autoequivalence
of $\mathrm{PSU}(3)_5$ extends by the identity on $\mathcal P$.
It is therefore trivial. Galois transport proves the same
statement for $\mathcal B$. Thus its nontrivial based duality
permutation determines a unique tensor autoequivalence class.
Moreover $U(\mathcal B)=1$, so
$\operatorname{Aut}_{\otimes}(\mathrm{Id}_{\mathcal B})=1$;
the coherent $C_2$ action realizing this class is unique up to
equivalence when it exists.

Finally, the first-column multiplicities give the actual
internal-end algebra
$\underline{\operatorname{End}}_{\mathcal E}(m_0)=1\oplus u$.
Its associativity trivializes the pointed associator of $u$.
The grading is consequently a crossed product
$\mathcal E=\mathcal B\rtimes C_2$, and Morita reflection
through this algebra identifies $\mathcal C$ with
$\mathcal B^{C_2}$.
The existing category in Proposition~\ref{prop:c4-centralizer}
has a central Tannakian boson, and its de-equivariantization
is a Galois conjugate of $\mathrm{PSU}(3)_5$ by
\cite[Section~2.4]{BRWRealizing}. Its action swaps exactly
the two dual pairs and is the same nontrivial based action.
The uniqueness above forces $K_0(\mathcal C)=R8\_2\_666$.
Every simple in that ring is self-dual, whereas labels $1,2$
in $R8\_2\_342$ form a non-self-dual pair. This is a
contradiction.
\end{proof}

\section{Near-group centralizer constructions and a common spherical reduction}
\label{sec:simple-galois-pu}

\begin{proposition}\label{prop:c4-centralizer}
The ring $R8\_2\_666$ admits a unitary braided categorification:
the centralizer of the full pointed subcategory in the center
of a near-group $C_4+4$ category.
\end{proposition}
\begin{proof}
The smallest hyperbolic case in
\cite[Section~3.2.3.1]{GIQuadraticData} is explicitly identified
with the actual centers of the two complex-conjugate near-group
$C_4+4$ categories. This finite case precedes the conjecture
there for larger groups. A convenient description is given in
\cite[Section~2.4]{BRWRealizing}:
\[
G=C_4^2,\quad q_G(x,y)=i^{-xy},\quad \theta_G(x,y)=(y,x),
\]
\[
H=C_{16}\times C_2,\quad
q_H(x,y)=\zeta_{32}^{3x^2}i^{y^2},\quad
\theta_H(x,y)=(3x+8y,x+y).
\]
The fixed groups identify with $K=C_4$ through
$k\mapsto(k,k)$ and $k\mapsto(4k,k\bmod2)$.
The center has $28$ simple objects, labelled by
$K$, $K\times\{\pi\}$ and representatives of the nonfixed
involution orbits in $G,H$.
Its pointed subcategory $\mathcal P$ has dimension four and
twist $q_K(k)=i^{-k^2}$. Trivial monodromy with every pointed
object selects exactly the eight labels
\[
(0,0),(2,2),(0,0,\pi),(2,2,\pi),
(3,1)_G,(2,1)_H,(6,1)_H,(4,0)_H.
\]
Thus $\mathcal P'$ is an existing unitary braided fusion
category, of dimension $96+64\sqrt2$. Its M\"uger center is
the Tannakian order-two radical of $\mathcal P$.

The source-bound certificate
\texttt{verify\_c4\_centralizer.py} reconstructs the full
$28$-label matrix from \cite[Equation~(1)]{BRWRealizing}.
It works in $\mathbb Z[\zeta_{32}]$, using
$\zeta_{32}^{16}=-1$ and
$\sqrt2=\zeta_{32}^4+\zeta_{32}^{-4}$.
It checks unitarity and all $1792$ Verlinde coefficients
whose first two labels belong to this centralizer, including
vanishing of every coefficient outside it.
The resulting $512$ coefficients identify its based ring
with $R8\_2\_666$ under the unital map
\[
(0,1,3,4,7,2,5,6).
\]
The proof note records the entire matrix formula, transversals,
normalization and source hash, and the normal and optimized
replays also verify associativity and self-duality.
\end{proof}

\begin{proposition}\label{prop:near-group-centralizers}
The rings $R8\_2\_820$ and $R8\_2\_869$ admit unitary modular
categorifications. They are the nonpointed M\"uger factors of
the Drinfeld centers of the first and second near-group
$C_5+5$ categories in \cite[Proposition~6 and Table~2]{EGNearGroup}.
\end{proposition}
\begin{proof}
Use the actual $C^*$-fusion categories with parameters
$(c,m)=(-1,1)$ and $(e^{2\pi i/3},2)$ in that table.
Their auxiliary metric groups are respectively
$G'=C_3^2$, $Q'(x,y)=(x^2+y^2)/3$, and
$G'=C_9$, $Q'(x)=2x^2/9$.
The center formulas of \cite[Proposition~7(b)]{EGNearGroup}
are identified with the actual centers for every odd group
order at most $13$ in the paragraph following its Conjecture~2.
Thus the order-five case here is established independently of
the conjecture for arbitrary groups.

Put
\[
a=\frac{5+3\sqrt5}{2},\qquad b=a+1,\qquad c=a+2.
\]
Each near-group category has global dimension $5c$.
The five invertible center objects $a_g$, $g\in C_5$, have
pairing $\exp(2\pi i\,2mgh/5)$, which is nondegenerate.
Their M\"uger centralizer is consequently an existing unitary
modular category of dimension $5c^2=45b$.
Trivial monodromy with these invertibles selects exactly
\[
a_0,\ b_0,\ c_{1,-1},\ c_{2,-2},\
d_{0,\gamma_1},\ldots,d_{0,\gamma_4},
\qquad
\{\gamma_i\}=(G'\setminus\{0\})/\{\pm1\}.
\]
Their dimensions are $(1,b,c,c,a,a,a,a)$; their squared
dimensions sum to $45b$, confirming completeness.

The exact source-bound certificate
\texttt{verify\_neargroup\_centralizer.py} reconstructs the
restricted $S$ matrix from those formulas and evaluates all
$512$ Verlinde coefficients over
$\mathbb Q[q,z]/(q^2-5,z^3-3z+1)$.
It identifies the two complete tensors under the respective
model-to-source maps
\[
(0,2,6,7,1,3,4,5),\qquad
(0,3,6,7,1,2,5,4).
\]
Here the auxiliary representatives are
$(1,0),(0,1),(1,1),(1,2)$ in $C_3^2$ and
$1,2,3,4$ in $C_9$.
An independent certificate instead reconstructs every tensor
entry from the integer fusion formulas following
\cite[Proposition~7(b)]{EGNearGroup} and obtains the same
identifications. Existence follows from the categories and
their centralizers; the finite computations identify their
based rings.
\end{proof}

The following result, obtained before these constructions,
also controls all possible categorifications of the two rings.
\begin{proposition}\label{prop:simple-galois-pu}
Every complex fusion category with based ring $R8\_2\_820$ or
$R8\_2\_869$ admits a spherical structure. Every spherical
dimension character belongs to the FP Galois orbit. Consequently,
categorifiability of either ring is equivalent to pseudo-unitary
categorifiability, and any categorification has a pseudo-unitary
Galois conjugate.
\end{proposition}
\begin{proof}
Both tensors are self-dual. In the canonical sphericalization
choose the positive-second-indicator lift of each simple. The
resulting signed quotient tensor $T$ is fully symmetric, has
$|T_{ij}^k|\leq N_{ij}^k$ and $T_{ij}^k\equiv N_{ij}^k\pmod2$,
and satisfies all associativity equations. The odd-channel parity
theorem \cite[Theorem~3.9]{TPEParity} gives the stronger restriction
\[
N_{jj}^i-T_{jj}^i\equiv0\pmod4.
\]
Every coefficient two in these tensors has a repeated label.
Thus repeated coefficients one are fixed at $+1$, repeated
coefficients two lie in $\{-2,2\}$, and the remaining
all-distinct coefficients one lie in $\{-1,1\}$.

The source-bound checker reconstructs exactly $35$ Boolean signs
for $R8\_2\_820$ and $36$ for $R8\_2\_869$. Expanding every
ordered associativity identity and reducing sign squares to one
leaves $215$ distinct integer polynomial equations for either
tensor. Exhaustive branching visits both signs at each unassigned
variable and rejects a branch only when a fully assigned equation
fails. It visits respectively $3071$ and $1873$ nodes and finds
exactly one solution in each case: all signs are positive.
Hence $T=N$. The chosen lifts form a full fusion subcategory of
the sphericalization, and the forgetful functor restricts to a
tensor equivalence with the original category. Restricting its
spherical structure proves the first assertion.

The remaining reduction is common and theoretical. Write the
basis as $(1,a_1,a_2,a_3,a_4,b,c_1,c_2)$, where $b$ has source
label $2$ for $R8\_2\_820$ and label $3$ for $R8\_2\_869$,
the $c$-labels are $6,7$, and the $a$-labels are the four other
nonunit labels. Both integral based rings satisfy
\[
b^2=1+\sum_{i=1}^4a_i+b+c_1+c_2,\qquad
(b+1)a_i=b^2-1,\qquad(b-1)c_j=b^2-1.
\]
Let $x=\dim(b)$ for any spherical structure. Nonzero simple
dimensions exclude $x=1$ and $x=-1$ by the last two identities.
They therefore give $\dim(a_i)=x-1$ and $\dim(c_j)=x+1$.
The first identity becomes
\[
x^2=1+4(x-1)+x+2(x+1)=7x-1,
\qquad x=\frac{7\pm3\sqrt5}{2}.
\]
These are exactly the two Galois-conjugate FP characters.
Conjugate the category together with its spherical structure to
the positive character. A spherical category with positive simple
dimensions is pseudo-unitary. This statement concerns every
categorification; it does not assert that each categorification
is unitarizable.
\end{proof}

\subsection{A unified character-parity exclusion for the Fibonacci pair}

\begin{proposition}\label{prop:fib-pair-galois-pu}
Every complex categorification of $R8\_2\_190$ or $R8\_2\_489$
admits a spherical structure and has a pseudo-unitary Galois
conjugate. For $R8\_2\_489$, the pointed $C_2$ subcategory has
trivial associator, so $1+g$ is a connected separable algebra.
The following theorem uses this structural reduction to exclude
both rings.
\end{proposition}
\begin{proof}
Both rings have basis $g^i\tau^jx^k$, $i,j,k\in\{0,1\}$, and
commutative presentation
\[
g^2=1,\qquad \tau^2=1+\tau,\qquad
x^2=g^\epsilon+\tau(x+gx),\qquad x^*=g^\epsilon x,
\]
where $\epsilon=1$ for $R8\_2\_190$ and $\epsilon=0$ for
$R8\_2\_489$. In the canonical sphericalization, choose the
lift of $g$ with dimension $1$. It squares to the unit.
A lift of $\tau$ is self-dual, and its signed square has the
form $1\pm\tau$; change the lift to make the sign positive.
Choose a lift of $x$, and define the other lifts by the simple
products $g\tau,gx,\tau x,g\tau x$. The relevant left and
right products have the same underlying simple and equal
nonzero spherical dimensions, so the chosen generators commute
in the signed quotient. The same argument gives
$x^*=g^\epsilon x$. The unit channel of $xx^*$ therefore fixes
the constant coefficient in
\[
x^2=g^\epsilon+a\tau x+b\tau gx,\qquad a,b\in\{1,-1\}.
\]
These relations determine the entire signed multiplication
tensor. When $\epsilon=1$, duality fixes $x^2$ and exchanges
the last two terms, so $a=b$.

For $\epsilon=0$, suppose instead $a=-b$. Writing
$t=\dim(\tau)$ and $z=\dim(x)$ gives $t^2=t+1$, $z^2=1$,
and global dimension $\delta=4(t+2)$.
A source character with values $g\mapsto1$, $\tau\mapsto t$,
$x\mapsto d$, where $d^2=1+2td$, has formal codegree
$f_E=2(t+2)(1+d^2)$. Central induction in the sphericalization
requires $\delta/f_E$ to be an algebraic integer
\cite[Theorem~2.13]{OstrikPiv3}: the sphericalization dimension
and the source formal codegree both acquire a factor two.
However,
\[
d^4-2d^3-6d^2+2d+1=0,\qquad
 y=\frac{\delta}{f_E}=\frac{2}{1+d^2}
\]
has primitive minimal polynomial
\[
5y^4-20y^3+25y^2-10y+1.
\]
Indeed the polynomial of $d$ is irreducible modulo three, and
$d=((d^2)^2-6d^2+1)/(2(d^2-1))$ gives
$\mathbb Q(y)=\mathbb Q(d)$; elimination gives the displayed
polynomial. Its leading coefficient excludes integrality.
Thus $a=b$ in both cases. Changing the lift of $x$ and its
three translates if necessary makes both signs positive.
Now $T=N$, so the chosen lifts form a tensor-closed spherical
section of the forgetful functor.

This section has $\dim(g)=1$, $t^2=t+1$, and $z^2=1+2tz$.
The irreducible quartic above shows that these four dimension
characters form a single Galois orbit, containing the positive
character
\[
t=\varphi=\frac{1+\sqrt5}{2},\qquad
z=\varphi+\sqrt{\varphi^2+1}.
\]
Conjugating to that character gives pseudo-unitarity.

For $R8\_2\_489$, work with this positive spherical structure.
A homogeneous pairing on a self-dual object $V$ induces one on
$V^n$ of sign $\nu_2(V)^n$. Any opposite-indicator self-dual
summand consequently has even multiplicity. Since
$\tau^2=1+\tau$, we have $\nu_2(\tau)=1$. The exact congruences
\[
x^3\equiv x+\tau+g\tau\pmod2,\qquad
(g\tau)^3\equiv g\pmod2
\]
then imply $\nu_2(x)=\nu_2(g\tau)=\nu_2(g)=1$.
For the positive spherical structure on
$\operatorname{Vec}_{C_2}^{\omega}$, the nontrivial associator
class has indicator $-1$ on $g$. Hence $\omega$ is trivial,
and $1+g$ is its regular separable algebra. This algebra also
exists before Galois conjugation. No half-braiding on $g$ is
inferred.

The source-bound replay
\texttt{verify\_fib\_pair\_theoretical.py} checks both complete
multiplication tensors, the duality and cube identities,
irreducibility modulo three, and the elimination polynomial.
The proof does not require an enumeration of signed tensors.
\end{proof}


\begin{theorem}\label{thm:fib-character-parity}
Neither $R8\_2\_190$ nor $R8\_2\_489$ admits a complex
fusion-category categorification.
\end{theorem}
\begin{proof}
By Proposition~\ref{prop:fib-pair-galois-pu}, it suffices to
exclude pseudo-unitary categorifications. Retain the common
presentation above, put $y=\tau x$, and write
\[
\varphi=\frac{1+\sqrt5}{2},\qquad
 d=\varphi+\sqrt{\varphi+2},\qquad
 D=2(\varphi+2)(1+d^2).
\]
The four dimension coordinates $1,\varphi,d,\varphi d$ are
rationally independent: the norm of $\varphi+2$ from
$\mathbb Q(\varphi)$ to $\mathbb Q$ is the nonsquare $5$,
so $d$ has degree two over $\mathbb Q(\varphi)$.

For a character $\chi$, let $Z_\chi$ denote its central simple
summand of $I(1)$ and let $\beta_\chi=D/f_\chi$.
The character correspondence preceding
\cite[Theorem~2.13]{OstrikPiv3}, together with its Example~2.9,
gives $\dim Z_\chi=\beta_\chi$, $[FZ_\chi:1]=1$, and
\[
S_{Z_\chi,W}=\beta_\chi\chi(FW)
\]
for the unnormalized modular matrix of the center. Symmetry
therefore yields
\begin{equation}\label{eq:fib-character-reciprocity}
\beta_\chi\chi(FZ_\psi)=\beta_\psi\psi(FZ_\chi).
\end{equation}

Take $\omega=1$ for $\epsilon=0$ and $\omega=i$ for
$\epsilon=1$. The two characters
\[
\chi_\pm(g)=-1,\qquad \chi_\pm(\tau)=\varphi,
\qquad \chi_\pm(x)=\pm\omega
\]
have formal codegree $4(\varphi+2)$ and central dimension
$\beta_\pm=1+\varphi d$. Formal codegrees here use
$\chi(b)\chi(b^*)$, which also gives this value when
$x^*=gx$ and $\omega=i$. Their central restrictions consequently
satisfy
\[
 FZ_\pm\in\{1+y,\,1+gy\},
\]
by the independence of the four dimension coordinates and unit
multiplicity one.

Now consider the character
\[
\psi(g)=1,\qquad \psi(\tau)=1-\varphi,
\qquad \psi(x)=z=(\varphi-1)d-\varphi.
\]
The relations give $z^2=1+2(1-\varphi)z$, and direct exact
reduction gives
\[
f_\psi=2(3-\varphi)(1+z^2),\qquad
\beta_\psi=(1+\varphi)+(2+3\varphi)d.
\]
Write $B=Z_\psi$. Because $\psi(g)=1$, the values
$\psi(FZ_+)$ and $\psi(FZ_-)$ coincide for every permitted
choice of the two restrictions. Applying
\eqref{eq:fib-character-reciprocity} twice and using
$\beta_+=\beta_-$ gives $\chi_+(FB)=\chi_-(FB)$.

The dimension of $B$ requires total multiplicity two for
$x,gx$ in $FB$ and total multiplicity three for $y,gy$.
If $\Delta_x,\Delta_y$ are the respective integer multiplicity
differences, subtracting these two character evaluations gives
\[
0=2\omega(\Delta_x+\varphi\Delta_y).
\]
Irrationality of $\varphi$ forces both differences to vanish.
The two equal multiplicities of $y$ and $gy$ cannot sum to three,
which is the required contradiction.

The source-bound replay
\texttt{verify\_fib\_character\_parity.py} checks all $512$
tensor entries for each presentation, the three characters,
the formal codegrees, and the exact dimension identities.
The obstruction itself requires no induction-column enumeration.
\end{proof}

\subsection{Cyclic induction phase multiplicities exclude $R8\_1\_9$}

\begin{lemma}\label{lem:cyclic-induction-phase-mass}
Let $a$ be an invertible object of order $n$ in a pseudo-unitary
fusion category of dimension $D$. Every phase occurring in $I(a)$
is an $n$th root of $\nu_n(a)$, and the total dimension in each of
these $n$ phases, counted with its induction multiplicity, is $D/n$.
\end{lemma}
\begin{proof}
The rotation spaces $\operatorname{Hom}(1,a^m)$ vanish for
$0<m<n$, so $\nu_m(a)=0$. For $m=n$ that space is
one-dimensional, and $\eta=\nu_n(a)$ is a root of unity.
The positive weighted trace
$\operatorname{Tr}(\theta^n I(a))=D\eta$ has maximal absolute
value $D=\dim I(a)$, so every constituent has $\theta^n=\eta$.
Choose $\lambda^n=\eta$ and write $W_j$ for the weight of phase
$\lambda\zeta_n^j$. The zeroth and first $n-1$ trace identities
are $\sum_jW_j=D$ and
$\sum_jW_j\lambda^m\zeta_n^{jm}=0$ for $1\leq m<n$.
Discrete Fourier inversion gives $W_j=D/n$.
\end{proof}

\begin{theorem}\label{thm:r9-center-trace}
The ring $R8\_1\_9$ has no complex fusion categorification.
\end{theorem}
\begin{proof}
By \cite[Lemma~3.1]{SchNear}, take a pseudo-unitary Galois conjugate
and put $d=2+\sqrt5$, $D=8+16d$. The odd-cube pivotal identities
force the second indicators, in source order, to be
$(1,0,0,0,0,1,1,1)$. The pointed generator $a$ has source label1,
with $a^{-1}$ of label2 and $a^2$ of label7. Its associator class
$k\in\mathbb Z/4$ satisfies $\nu_2(a^2)=(-1)^k$, so $k=0$ or2.
The canonical pointed indicators give
$\nu_4(a)=\nu_4(a^{-1})=(-1)^{k/2}$.

The pointed induction Gram block is $8I_4$. Hence a central row
$v$ occurring over $a$ has no other pointed coordinate, and its
dimension has the form $v_a+Bd$, with $B\geq0$.
Lemma~\ref{lem:cyclic-induction-phase-mass} gives weight
$D/4=2+4d$ in each actual phase. Comparing constant coefficients
therefore gives $\sum_{Z\text{ in one phase}}v_a(Z)^2=2$.
Thus $v_1,v_2\leq1$ for every central row. This argument is uniform
in the pointed associator class and does not average away the two
imaginary phases.

The complete formal codegrees are $D,D',8,8,8,8,8,8$;
the six other characters lie in $\mathbb Q(i)$. Actual
character-labelled Galois symmetry transports a forgetful row by
\[
 T(v)=v+(0,0,0,1,1,1,1,0)
 \left(v_0+v_1+v_2+v_7-\frac{v_3+v_4+v_5+v_6}{2}\right),
\]
up to the sign making it nonnegative. Partner rows have equal
normalized relative moments, because their twists are related by
$\sigma^2$ and $\sigma^2$ fixes $\mathbb Q(\sqrt5)$.
The two forced $I(1)$ rows and the six remaining
character-labelled rows leave exactly eight possible low frames.

For arbitrary roots of unity, the first two relative traces lie
in a fixed $17$-vertex hull. It consists of the nine pairs from
orders $1,2,4,5,10,20$, together with the needed negative-half
endpoints: an extra odd factor coprime to $10$ scales both moments
by $\mu(m)/\varphi(m)\in[-1/2,1]$. No root-order bound or
prime-support assumption is used. The full elementary Gauss-sum
proof is supplied in the evidence packet.

For $k=0$, constituents over labels1 and2 have $\theta^4=1$;
for $k=2$ they have $\theta^4=-1$ and their first two relative
traces vanish. Those over label7 have $\theta^2=1$.
Use these phase restrictions, the exact dimension-integrality
condition $D/\dim(Z)$ integral, Galois pairing, and the residual
Gram bound for the entire row orbit. With the Fourier bound,
class zero leaves $64$ row orbits and $405$ moment columns per
frame. The original, weaker bounds suffice for class two and leave
$92$ orbits and $389$ columns.

The Gram and two induction trace identities have $56$ integer
coordinates. The two covectors in
\texttt{r9\_center\_trace\_certificate.json} have nonnegative
pairing with every column, but target pairings $-64$ in class zero
and $-6432$ in class two, for all eight low frames. Their standalone
standard-library verifier reconstructs every bounded domain and
every positive semidefinite condition with exact rational
arithmetic. Neither target is a nonnegative combination of its
columns, excluding both possible pointed classes.
The initial Galois reduction proves the general complex assertion.
\end{proof}

\subsection{The full modular identity and a fourth-moment projection}

\begin{lemma}\label{lem:full-induction-sts}
Let $\mathcal C$ be a pseudo-unitary fusion category with commutative
Grothendieck ring, and let $b_\chi$ be the forgetful vector of the
actual summand $A_\chi\subset I(1)$ corresponding to the character
$\chi$. For a simple $Z\in Z(\mathcal C)$ write $v_Z$ for its
forgetful vector. Then
\begin{equation}\label{eq:full-induction-sts}
 \sum_Z\theta_Z(v_Z)_i(v_Z)_j
 =\sum_\chi(b_\chi)_i\overline{\chi(X_j)}.
\end{equation}
\end{lemma}
\begin{proof}
Let $C_{\chi i}=\chi(X_i)$, let $f_\chi$ be the formal codegrees,
and put $X=\operatorname{diag}(f_\chi^{-1})C$.
The character identification in \cite[Example~2.9 and
Theorem~2.13]{OstrikPiv3} gives
$S_{A_\chi,Z}=\chi(FZ)/f_\chi$ for every central simple $Z$,
with $S_{1Z}=\dim(Z)/\dim(\mathcal C)$.
All $A_\chi$ have twist one, and the center has trivial anomaly.
Thus the low block of $STS=T^{-1}ST^{-1}$ is
\[
 X\left(\sum_Z\theta_Zv_Zv_Z^{\mathsf t}\right)X^{\mathsf t}
 =XB^{\mathsf t},
\]
where $B$ has rows $b_\chi$. Character orthogonality gives
$C\overline C^{\mathsf t}=\operatorname{diag}(f_\chi)$,
so $X^{-\mathsf t}=\overline C$. Cancellation proves the identity.
\end{proof}

\begin{theorem}\label{thm:r18-sts-moments}
The ring $R8\_1\_18$ has no complex fusion categorification.
\end{theorem}
\begin{proof}
By \cite[Lemma~3.1]{SchNear}, pass to a pseudo-unitary Galois
conjugate. Put $d=2+\sqrt5$, $K=\mathbb Q(\sqrt5)$ and
$D=8+16d$. In model order
$(1,a,a^2,a^3,\rho,a\rho,a^2\rho,a^3\rho)$, mapped to source
$(0,1,7,2,3,6,5,4)$, the ring is commutative and satisfies
\[
 a^4=1,\qquad \rho^2=a+(1+a+a^2+a^3)\rho.
\]
Its six non-FP characters have formal codegree eight and are
\[
 \chi_{k,\pm}(a)=\zeta_8^{2k},\qquad
 \chi_{k,\pm}(\rho)=\pm\zeta_8^k,\qquad k=1,2,3.
\]
The other two codegrees are $D,D'$.
The two forced low rows are $1$ and $1+X_3+X_4+X_5+X_6$;
the other six have form $1+$ two nonpointed constituents.
Actual character reciprocity and duality leave eight low assignments.
Taking normalized relative trace $t_K$ in
\eqref{eq:full-induction-sts} gives, for nonunit indices,
\[
 \sum_Z(v_Z)_i(v_Z)_j t_K(\theta_Z)
 =\begin{cases}4-d&i,j\in\{3,4,5,6\},\\0&\text{otherwise}.
 \end{cases}
\]
Indeed every nonpointed value of the six $\mathbb Q(\zeta_8)$
characters has relative trace zero; the pointed entries cancel in
each of the eight assignments.

The pointed induction Gram is $8I_4$.
Lemma~\ref{lem:cyclic-induction-phase-mass} therefore gives
$(v_Z)_1,(v_Z)_2\leq1$ exactly as in the preceding theorem,
uniformly over all four pointed associator classes.
Besides the quadratic transport $T$ displayed there, full Galois
symmetry acts on rows by the source permutation
\[
 P=(0,2,1,6,5,4,3,7).
\]
It is induced by an automorphism fixing $K$ and sending
$\zeta_8$ to $\zeta_8^3$: on the model basis it sends
$a\mapsto a^{-1}$ and $a^j\rho\mapsto a^{3j+1}\rho$.
Include also source duality
$(0,2,1,4,3,6,5,7)$. Every full row orbit has a common
relative moment histogram: Galois transport changes a twist by
$\sigma^2$, which fixes quadratic $K$, and duality preserves twists.

The Cauchy theorem for spherical fusion categories
\cite[Theorem~3.9]{BNRWCauchy} applies to the nonintegral
dimension $D$. Its norm is $320=2^6\cdot5$, so the exponent has
only primes two and five. For powers $1,2,4$, every normalized
relative twist triple is one of the triples from fortieth roots
or the zero triple. This follows directly from the quadratic
Gauss sums: fifth-adic exponent at least two, or $2$-adic exponent
at least four, gives zero in all three coordinates. There are
exactly thirteen triples; this is not an assertion that the
exponent divides forty.

There is no need to determine fourth indicators. Each $\nu_4(X_i)$
is a sum of fourth roots, so $t_K(\nu_4(X_i))\in\mathbb Z$.
The functional $\ell(a+bd)=2a-b$ annihilates $D\mathbb Q$.
Applying it to the fourth induction identity gives
\[
 \ell\!\left(\sum_Z(v_Z)_i\dim(Z)t_K(\theta_Z^4)\right)=0.
\]
Together with the first two induction traces, the full modular
identity above, and the Gram identity, this gives $119$ rational
linear coordinates after low terms are subtracted.

The exact domain uses the diagonal Gram bounds, integrality of
$D/\dim(Z)$, the Fourier multiplicity bound, full row-orbit
closure, and positive semidefiniteness after subtracting that
entire orbit. It has $25$ row orbits for every low assignment.
The pointed classes $k=0$, $k=1,3$ and $k=2$ impose their respective
cyclic twist constraints and leave $167$, $145$ and $153$ columns.
The three integer covectors in
\texttt{r18\_sts\_moment\_certificate.json} pair nonnegatively
with every column and pair to $-256$ with every corresponding
target, for all eight low assignments. The standard-library
verifier reconstructs the source, characters, thirteen twist
triples and all bounded domains using exact rational arithmetic.
Thus none of the four pointed classes is possible. The initial
Galois reduction proves the general complex assertion.
\end{proof}


\subsection{A local obstruction to an inner square root}

\begin{lemma}\label{lem:2d2-neutral-indicators}
Let $\mathcal B$ be pseudo-unitary with simple objects $1,g,r,gr$
and rules
\[
 g^2=1,\qquad gr=rg,\qquad r^2=1+2r+2gr.
\]
Then $\nu_2(g)=-1$, $\nu_2(r)=\nu_2(gr)=1$, and the cubic
rotations on $\operatorname{Hom}(1,r^3)$ and
$\operatorname{Hom}(1,(gr)^3)$ are identities.
\end{lemma}
\begin{proof}
Put $d=2+\sqrt5$ and $D=4+8d=20+8\sqrt5$.
The formal codegrees are $D,D',4,4$, and
$1/D-1/D'=-1/\sqrt5$. The proof of Lemma~\ref{lem:q17-conductor-equality}, with the
primitive character modulo five in place of that modulo seventeen,
gives the same equality conclusion at conductor five: every central simple with nonrational dimension
ratio $z/D$ has twist of order dividing five.

We use the signed version of the saturated indicator identity.
Let $\ell(z)$ be a rational coordinate of the nontrivial quadratic
part of $z/D$, scaled so that its values on $(1,g,r,gr)$ are
$(2,2,-1,-1)$. For a cyclotomic automorphism $\sigma$ with residue
$k$ modulo five, write $\pi$ and $\varepsilon_A$ for its modular
permutation and signs. The Chinese remainder theorem lifts each chosen
unit residue to an exponent coprime to the full cyclotomic conductor.
The dimension transformation gives
\[
 z_{\pi A}=\varepsilon_A\frac D{\sigma D}\sigma(z_A),
 \qquad \ell(z_{\pi A})=\varepsilon_A\chi_5(k)\ell(z_A).
\]
The first induction identity, weighted by $\ell$, is
\[
 \sum_A\ell(z_A)z_A\theta_A=D\ell(1).
\]
The signs in the two preceding transformations cancel, giving
\[
 \sigma\bigl(\ell(z_A)z_A\bigr)
 =\chi_5(k)\frac{\sigma D}{D}\ell(z_{\pi A})z_{\pi A}.
\]
Apply $\sigma$ to the weighted identity and reindex by $\pi$.
If $h$ is the inverse of $k$ modulo five, its nonzero terms obey
$\sigma(\theta_A)=\theta_{\pi A}^{h}$, since their twists divide
five and $\theta_{\pi A}=\sigma^2(\theta_A)$. The terms belonging
to $I(1)$ have twist one and satisfy the same equation.
Cancel $\sigma D$ and use the induction formula for the $h$th
indicator. Inversion permutes the units modulo five and
$\chi_5(h)=\chi_5(k)$, proving
\[
 \sum_i\ell(\dim X_i)\nu_k(X_i)=\chi_5(k)\ell(1)
 \qquad (k,5)=1.
\]
At $k=2$ this reads
$2+2\nu_2(g)-\nu_2(r)-\nu_2(gr)=-2$.
All three unknowns are signs, forcing the stated values.
At $k=3$ it reads $\nu_3(r)+\nu_3(gr)=4$.
Each rotation acts on a two-dimensional space and has order
three. Equality in the trace bound forces both rotations to be
identities.
\end{proof}

\begin{lemma}\label{lem:2d2-no-inner-square}
For such a pseudo-unitary $\mathcal B$, no tensor autoequivalence
fixing all four simple labels has square tensor-isomorphic to
$\operatorname{Ad}(g)$.
\end{lemma}
\begin{proof}
We use only the indicated local pivotal identities from
\cite[\S3.1.1 and pp.~539--540]{EIP}, interpreted on opposite
composition-dual Hom spaces using \cite[Lemma~3.13(ii),
Proposition~3.22 and Corollary~4.9]{FGPartialDual}.
No Hermitian form or unitarity assumption is made.
Put
\[
 H^+=\operatorname{Hom}(r^2,r),\quad
 J^+=\operatorname{Hom}(r^2,gr),\qquad
 H^-=\operatorname{Hom}(r,r^2),\quad
 J^-=\operatorname{Hom}(gr,r^2).
\]
Each space has dimension two. Normalize the $g$--$r$ crossing
$c$ so that its pivotal one-click rotation is its inverse.
Its bend parameter is $\mu=\nu_2(gr)/\nu_2(r)=1$.

Here is the paired-space normalization, including the distinction
between transpose and inverse transpose. For $F=\operatorname{Ad}(g)$,
take
$q_r=(1_r\otimes\operatorname{ev}_g)(c\otimes1_g)$ and
$q_{gr}=1_g\otimes q_r$. On opposite Hom spaces its outgoing and
incoming actions use respectively these identifications and their
inverses. Suppressing the tensorator and its inverse, they are
\[
 T^+(f)=q_BF(f)(q_A^{-1}\otimes q_A^{-1}),\qquad
 T^-(h)=(q_A\otimes q_A)F(h)q_B^{-1}.
\]
Composition pairing is invariant, so the incoming loop
matrix is $K^-=(K^+)^{-\mathsf t}$. Transposes of the object maps
are relative to the canonical pivotal structure of $F$; the unit
map $q_1=\operatorname{ev}_g$ has form defect one.
The two defects are $s_r=\mu$ and $s_{gr}=\mu\nu_2(g)$:
the first contraction is the defining two-crossing $\mu$ diagram,
and the second reverses one further antisymmetric $g$ form.
For a right bend $R$ on $\operatorname{Hom}(A^2,B)$ this gives
$RT^+=(s_B/s_A^2)T^-R$.

The direct loop contractions give $K_H^2=-1$, $K_J^2=1$.
The nondegenerate symmetric right theta pairings consequently pair
opposite loop eigenlines. We may choose composition-dual bases with
\[
 \begin{gathered}
 K_H^+=\operatorname{diag}(i,-i),\quad
 K_H^-=\operatorname{diag}(-i,i),\\
 K_J^+=K_J^-=\operatorname{diag}(1,-1),\qquad
 R_H^+=R_H^-=R_J^+=R_J^-=E=\begin{pmatrix}0&1\\1&0\end{pmatrix}.
 \end{gathered}
\]
Their remaining basis changes are
$\operatorname{diag}(u,u^{-1})$ and
$\operatorname{diag}(v,v^{-1})$.
The cubic rotation on $H$ is the identity by
Lemma~\ref{lem:2d2-neutral-indicators}. The two decorated cubic
phases on $J$ remain arbitrary third roots throughout the argument.

In these bases the tetrahedral slots of \cite[pp.~528--530]{EIP}
have the required forms
\[
 A=\begin{pmatrix}
 t&0&0&A_1\\0&t&t&0\\0&t&t&0\\a_1&0&0&t
 \end{pmatrix},\qquad
 D_0=\begin{pmatrix}
 x&0&0&Z\\0&X&y&0\\0&Y&-X&0\\z&0&0&-x
 \end{pmatrix}.
\]
The odd partial-dual slot map is
$(i,j,k,l)\mapsto(k,1-j,i,1-l)$, pairing
$a_1\leftrightarrow A_1$, $x\leftrightarrow X$,
$y\leftrightarrow Y$ and $z\leftrightarrow Z$.
These are independent dual coefficients, not complex conjugates.
For completeness, the even slot maps are
\[
 c(i,j,k,l)=(j,1-k,1-i,l),\qquad
 q(i,j,k,l)=(1-k,1-l,1-i,1-j).
\]
The $A$ coefficients obey both even symmetries. For $D_0$, only
its eight odd-parity zeros and the two relations
$(D_0)_{1111}=-(D_0)_{0000}$ and
$(D_0)_{1010}=-(D_0)_{0101}$ are needed. In both flips $i=k$, so the relevant decorated phase
cancels as its own cube. No equality of decorated phases is used.

The fusion-ring corner on the four-dimensional annular module
over $r$ has block matrices
\[
 P=\begin{pmatrix}0&I\\I&0\end{pmatrix},\qquad
 Q=\begin{pmatrix}\phi&\phi_0\\\phi_0&\phi\end{pmatrix},
 \qquad \phi=2tI.
\]
The FP character is absent, since its associated central simple
is the central unit, whose forgetful image contains no $r$
\cite[\S2.3]{OstrikPiv3}. Thus $\phi+\phi_0=(2-\sqrt5)I$.
On $P=-1$ the fusion relation is $Q^2=I$, giving
\[
 t=\frac{2-\sqrt5+\tau}{4},\qquad \tau=\pm1.
\]
Three diagonal composition-dual resolutions on \cite[p.~539]{EIP}
now suffice:
\[
 t^2+xX=\frac{3-\sqrt5}{2},\qquad
 xX+zZ=\frac12,\qquad xX+yY=\frac12.
\]
They come from the first resolution at $(0,0,0,0)$ and the third
at $(0,0,0,0)$ and $(0,1,0,1)$. In the first identity the
unit-channel contribution is $1/d=\sqrt5-2$, so its uncombined
form is $1=2t^2+2xX+(\sqrt5-2)$. Each diagonal contraction pairs
every vertex phase with its reciprocal. Therefore
\[
 yY=zZ=
 \begin{cases}(\sqrt5-1)/8&\tau=1,\\
 (3\sqrt5-5)/8&\tau=-1.
 \end{cases}
\]
Both are nonzero, so $y,z\ne0$ in every branch.

Let $T$ fix the four simple labels. It preserves the canonical
pivotal structure. Normalize its object identifications to preserve
the self-dual forms. The crossing is taken to $\lambda c$ with
$\lambda=\pm1$, and its $H$ action satisfies
$T_HK_H=\lambda K_HT_H$. If $\lambda=-1$, its matrix is
antidiagonal, so $T_H^2$ is scalar. If $\lambda=1$, the normalized
crossing and hence both loop eigenspace pairs are fixed. The right
pairings then give diagonal actions
$\operatorname{diag}(u,u^{-1})$ and
$\operatorname{diag}(v,v^{-1})$. The two nonzero coefficients
transform with weights $u^{-2}v^2$ and $u^{-2}v^{-2}$.
Their invariance forces $u^4=1$, so again $T_H^2$ is scalar.

The projective action on $H$ depends only on the tensor-isomorphism
class of the autoequivalence and respects composition: changing
the identification of a simple object changes this action only
by a scalar. But $\operatorname{Ad}(g)$ acts projectively by
$K_H=\operatorname{diag}(i,-i)$, which is not scalar.
It therefore cannot be the square of $T$.
\end{proof}

\begin{theorem}\label{thm:r24-inner-square}
The ring $R8\_2\_24$ has no complex fusion categorification.
\end{theorem}
\begin{proof}
The source has exactly two FP dimensions, $1$ and $2+\sqrt5$,
so pass to a pseudo-unitary Galois conjugate using
\cite[Lemma~3.1]{SchNear}. The source tensor has a faithful
$C_2$ grading with neutral labels $\{0,5,6,7\}$ and the ring of
Lemma~\ref{lem:2d2-neutral-indicators}, with $r=X_5$ and $g=X_7$.
The odd invertible $a=X_1$ has $a^2=g$. Since the source ring is
commutative, $\operatorname{Ad}(a)$ fixes every neutral simple
label. Its square is tensor-isomorphic to
$\operatorname{Ad}(g)$, contradicting
Lemma~\ref{lem:2d2-no-inner-square}.
\end{proof}

The source-bound replay in
\texttt{verify\_r24\_projective\_inner\_square.py} checks
the neutral tensor, complete four-character table, saturation arithmetic,
coefficient-slot substitutions and both nonzero scalar branches.
The categorical argument above supplies the obstruction; the finite
replay does not assume a classification of all $2D2$ associators.

